\documentclass[a4paper]{article}

\usepackage{amsmath, amssymb, amsthm}
\usepackage{hyperref}
\usepackage[capitalize,nameinlink]{cleveref}
\usepackage{blueprint}

\newcommand{\C}{\mathbb{C}}
\newcommand{\eps}{\varepsilon}
\newcommand{\dist}{\operatorname{dist}}

\title{Blueprint for Runge's Theorem, Polynomial Form}
\author{}
\date{}

\begin{document}

\maketitle

\begin{abstract}
This blueprint proves the connected-complement polynomial form of Runge's
theorem: if \(K \subset U \subset \C\), where \(U\) is open, \(K\) is compact,
\(\C \setminus K\) is connected, and \(f\) is holomorphic on \(U\), then \(f\)
can be approximated uniformly on \(K\) by polynomials.
\end{abstract}

\tableofcontents

\section{Main statement}

\begin{theorem}[Runge's theorem, polynomial form]
\label{thm:main}
\lean{MainTheorem}
Let \(U,K \subset \C\). Assume that \(U\) is open, \(K\) is compact,
\(K \subset U\), \(\C \setminus K\) is connected, and \(f : \C \to \C\) is
holomorphic on \(U\). Then for every \(\eps > 0\), there exists a polynomial
\(p \in \C[X]\) such that
\[
  \forall z \in K,\qquad |p(z)-f(z)| < \eps .
\]
Equivalently, in Lean:
\[
\begin{aligned}
&\texttt{theorem MainTheorem \{U K : Set ℂ\} \{f : ℂ → ℂ\}}\\
&\quad \texttt{(hU : IsOpen U) (hK : IsCompact K) (hKU : K ⊆ U)}\\
&\quad \texttt{(hKc : IsConnected (Kᶜ)) (hf : DifferentiableOn ℂ f U) :}\\
&\quad \texttt{∀ ε > 0, ∃ p : Polynomial ℂ, ∀ z ∈ K, ‖p.eval z - f z‖ < ε.}
\end{aligned}
\]
\end{theorem}

\begin{proof}
Fix \(\eps > 0\). The proof has two main steps.

First, by a rational approximation version of Runge's theorem, since \(f\) is
holomorphic on an open neighbourhood \(U\) of the compact set \(K\), there is a
rational function \(r\), with all poles outside \(K\), such that
\[
  \forall z \in K,\qquad |r(z)-f(z)| < \eps/2 .
\]
More concretely, \(r\) may be taken to be a finite sum of simple poles:
\[
  r(z) = \sum_{j=1}^N \frac{c_j}{z-a_j},
  \qquad a_j \notin K .
\]

Second, because \(\C \setminus K\) is connected, each pole \(a_j\) can be
``moved to infinity'' inside \(\C \setminus K\). Along this path, the local
geometric-series identity allows one to approximate
\[
  \frac{1}{z-a_j}
\]
uniformly on \(K\) by a rational function whose pole is farther away. Iterating
this along a finite subdivision of the path moves the pole to a point \(b_j\)
with very large modulus. Once the pole is sufficiently far away, the expansion
\[
  \frac{1}{z-b_j}
  =
  -\frac{1}{b_j}\frac{1}{1-z/b_j}
  =
  -\sum_{n=0}^{\infty} \frac{z^n}{b_j^{n+1}}
\]
converges uniformly on \(K\). Hence \((z-b_j)^{-1}\), and therefore also
\((z-a_j)^{-1}\), can be uniformly approximated on \(K\) by polynomials.

Applying this to the finitely many poles \(a_1,\dots,a_N\), with error bounds
chosen small enough so that the total error is \(<\eps/2\), we obtain a
polynomial \(p\) such that
\[
  \forall z \in K,\qquad |p(z)-r(z)| < \eps/2 .
\]
Therefore, for all \(z \in K\),
\[
  |p(z)-f(z)|
  \le |p(z)-r(z)| + |r(z)-f(z)|
  < \eps/2+\eps/2
  = \eps .
\]
This proves the theorem.
\end{proof}

\section{Analytic preparation}

\begin{lemma}[Compact subset of an open set has a positive collar]
\label{lem:compact_subset_open_collar}
\lean{compact_subset_open_has_thickening}
Let \(K,U \subset \C\). If \(K\) is compact, \(U\) is open, and \(K \subset U\),
then there exists \(\rho > 0\) such that
\[
  \{z \in \C : \dist(z,K) < \rho\} \subset U .
\]
\end{lemma}

\begin{proof}
Since \(K\) is compact and \(U^c\) is closed, and since \(K \cap U^c=\varnothing\),
the distance between \(K\) and \(U^c\) is positive. Choose \(\rho\) smaller than
this distance.
\end{proof}

\begin{lemma}[Cauchy integral approximation by finite pole sums]
\label{lem:cauchy_integral_to_pole_sum}
\lean{cauchy_integral_approximated_by_pole_sum}
Let \(K \subset U \subset \C\), with \(K\) compact and \(U\) open. Let
\(f : \C \to \C\) be holomorphic on \(U\). Then for every \(\eps > 0\), there
exist finitely many points \(a_1,\dots,a_N \in \C \setminus K\) and coefficients
\(c_1,\dots,c_N \in \C\) such that
\[
  \forall z \in K,\qquad
  \left|
    \sum_{j=1}^N \frac{c_j}{z-a_j} - f(z)
  \right| < \eps .
\]
\end{lemma}

\begin{proof}
Using \cref{lem:compact_subset_open_collar}, choose a small neighbourhood of
\(K\) whose closure is contained in \(U\). Cover \(K\) by finitely many small
closed squares whose union is contained in this neighbourhood. Let \(\Gamma\)
be the positively oriented boundary of the union of these squares, after
cancelling internal edges.

For \(z \in K\), Cauchy's integral formula gives
\[
  f(z)
  =
  \frac{1}{2\pi i}
  \int_\Gamma \frac{f(\zeta)}{\zeta-z}\,d\zeta .
\]
The curve \(\Gamma\) is compact and disjoint from \(K\), so the kernel
\[
  (\zeta,z) \mapsto \frac{f(\zeta)}{\zeta-z}
\]
is uniformly continuous on \(\Gamma \times K\). Therefore the above contour
integral can be uniformly approximated, for \(z \in K\), by Riemann sums of the
form
\[
  \sum_{j=1}^N \frac{c_j}{a_j-z}.
\]
Changing signs in the coefficients gives the desired form
\[
  \sum_{j=1}^N \frac{c_j}{z-a_j}.
\]
Each pole \(a_j\) lies on \(\Gamma\), hence outside \(K\).
\end{proof}

\begin{proposition}[Rational approximation with poles off \(K\)]
\label{prop:rational_approximation}
\lean{rational_approximation_with_poles_off}
Let \(K \subset U \subset \C\), with \(K\) compact and \(U\) open. Let
\(f : \C \to \C\) be holomorphic on \(U\). Then for every \(\eps > 0\), there
exist \(N : \mathbb{N}\), coefficients \(c_1,\dots,c_N \in \C\), and poles
\(a_1,\dots,a_N \in \C \setminus K\), such that
\[
  \forall z \in K,\qquad
  \left|
    \sum_{j=1}^N \frac{c_j}{z-a_j} - f(z)
  \right| < \eps .
\]
\end{proposition}

\begin{proof}
This is exactly \cref{lem:cauchy_integral_to_pole_sum}.
\end{proof}

\section{Moving poles to infinity}

\begin{lemma}[The complement of a compact set is unbounded]
\label{lem:compact_complement_unbounded}
\lean{compact_complement_unbounded}
Let \(K \subset \C\) be compact. Then for every \(R > 0\), there exists
\(b \in \C \setminus K\) such that \(|b| > R\).
\end{lemma}

\begin{proof}
A compact subset of \(\C\) is bounded. Choose \(b\) with norm larger than both
\(R\) and a bound for \(K\).
\end{proof}

\begin{lemma}[Connected open subsets of \(\C\) are path-connected]
\label{lem:connected_open_path_connected}
\lean{isConnected_isOpen_pathConnected_complex}
Let \(V \subset \C\). If \(V\) is open and connected, then \(V\) is
path-connected.
\end{lemma}

\begin{proof}
This is the standard theorem that open connected subsets of \(\mathbb{R}^n\),
hence of \(\C\), are path-connected. One proves that the path component of a
point is both open and closed in \(V\), and then uses connectedness.
\end{proof}

\begin{lemma}[Path to infinity in the connected complement]
\label{lem:path_to_infinity}
\lean{path_to_infinity_in_connected_complement}
Let \(K \subset \C\) be compact and assume \(\C \setminus K\) is connected. If
\(a \notin K\), then for every \(R > 0\), there exist \(b \notin K\) and a path
\(\gamma : [0,1] \to \C \setminus K\) such that
\[
  \gamma(0)=a,\qquad \gamma(1)=b,\qquad |b|>R .
\]
\end{lemma}

\begin{proof}
The complement \(\C \setminus K\) is open because \(K\) is compact, hence
closed. By \cref{lem:compact_complement_unbounded}, choose \(b \in \C \setminus K\)
with \(|b|>R\). Since \(\C \setminus K\) is open and connected, it is
path-connected by \cref{lem:connected_open_path_connected}. Hence there is a
path from \(a\) to \(b\) inside \(\C \setminus K\).
\end{proof}

\begin{lemma}[Local pole-moving lemma]
\label{lem:local_pole_moving}
\lean{local_pole_moving}
Let \(K \subset \C\) be compact. Let \(a,b \in \C \setminus K\). Suppose that
\[
  |a-b| < \dist(b,K).
\]
Then for every \(\eps > 0\), there exists \(N : \mathbb{N}\) such that
\[
  \forall z \in K,\qquad
  \left|
    \frac{1}{z-a}
    -
    \sum_{n=0}^{N}
      \frac{(a-b)^n}{(z-b)^{n+1}}
  \right|
  < \eps .
\]
Thus the pole at \(a\) can be uniformly approximated on \(K\) by a rational
function whose only pole is at \(b\).
\end{lemma}

\begin{proof}
For \(z \in K\),
\[
  \frac{1}{z-a}
  =
  \frac{1}{z-b-(a-b)}
  =
  \frac{1}{z-b}
  \frac{1}{1-\frac{a-b}{z-b}} .
\]
Since \(|a-b| < \dist(b,K) \le |z-b|\) for all \(z \in K\), the geometric
series converges uniformly on \(K\). Taking a sufficiently long partial sum
gives the desired uniform error bound.
\end{proof}

\begin{lemma}[Pole sufficiently far away is polynomially approximable]
\label{lem:far_pole_polynomial}
\lean{far_pole_polynomial_approx}
Let \(K \subset \C\) be compact. Let \(b \in \C\) be such that
\[
  |z| < |b|
  \qquad \text{for all } z \in K .
\]
Then for every \(\eps > 0\), there exists a polynomial \(p \in \C[X]\) such that
\[
  \forall z \in K,\qquad
  \left|p(z)-\frac{1}{z-b}\right| < \eps .
\]
\end{lemma}

\begin{proof}
For \(z \in K\),
\[
  \frac{1}{z-b}
  =
  -\frac{1}{b}\frac{1}{1-z/b}
  =
  -\sum_{n=0}^{\infty} \frac{z^n}{b^{n+1}} .
\]
The series converges uniformly on \(K\) because \(|z/b|\) is uniformly bounded
by a constant \(<1\). A sufficiently long partial sum is a polynomial and gives
the desired approximation.
\end{proof}

\begin{proposition}[Single pole is polynomially approximable]
\label{prop:single_pole_polynomial}
\lean{single_pole_polynomial_approx}
Let \(K \subset \C\) be compact and assume \(\C \setminus K\) is connected. If
\(a \notin K\), then for every \(\eps > 0\), there exists a polynomial
\(p \in \C[X]\) such that
\[
  \forall z \in K,\qquad
  \left|p(z)-\frac{1}{z-a}\right| < \eps .
\]
\end{proposition}

\begin{proof}
Choose \(R\) so large that every \(z \in K\) satisfies \(|z| < R\). By
\cref{lem:path_to_infinity}, choose a point \(b \in \C \setminus K\) with
\(|b|>R\), and a path from \(a\) to \(b\) inside \(\C \setminus K\).

The image of this path is compact and contained in the open set
\(\C \setminus K\). Hence it has positive distance from \(K\). By uniform
continuity of the path, subdivide it into finitely many points
\[
  a=a_0,a_1,\dots,a_m=b
\]
such that
\[
  |a_i-a_{i+1}| < \dist(a_{i+1},K)
\]
for every \(i\). Applying \cref{lem:local_pole_moving} repeatedly, the function
\((z-a)^{-1}\) is uniformly approximated on \(K\) by a rational function whose
only pole is \(b\). Since \(|b|\) is larger than \(|z|\) for all \(z \in K\),
\cref{lem:far_pole_polynomial} approximates this final pole by a polynomial.
Combining the finitely many error estimates gives the desired polynomial
approximation.
\end{proof}

\begin{proposition}[Finite pole sums are polynomially approximable]
\label{prop:finite_pole_sum_polynomial}
\lean{finite_pole_sum_polynomial_approx}
Let \(K \subset \C\) be compact and assume \(\C \setminus K\) is connected. Let
\(a_1,\dots,a_N \in \C \setminus K\) and \(c_1,\dots,c_N \in \C\). Then for
every \(\eps > 0\), there exists a polynomial \(p \in \C[X]\) such that
\[
  \forall z \in K,\qquad
  \left|
    p(z) -
    \sum_{j=1}^N \frac{c_j}{z-a_j}
  \right|
  < \eps .
\]
\end{proposition}

\begin{proof}
For each \(j\), apply \cref{prop:single_pole_polynomial} to approximate
\((z-a_j)^{-1}\) by a polynomial \(q_j\). Choose the error for \(q_j\) small
enough so that after multiplication by \(|c_j|\) and summation over the finite
set \(j=1,\dots,N\), the total error is \(<\eps\). Then
\[
  p(z) := \sum_{j=1}^N c_j q_j(z)
\]
is a polynomial and satisfies the required bound.
\end{proof}

\section{Final assembly}

\begin{proof}[Proof of \cref{thm:main}]
Let \(\eps > 0\). By \cref{prop:rational_approximation}, choose a finite pole
sum
\[
  r(z) = \sum_{j=1}^N \frac{c_j}{z-a_j},
  \qquad a_j \notin K,
\]
such that
\[
  \forall z \in K,\qquad |r(z)-f(z)| < \eps/2 .
\]
By \cref{prop:finite_pole_sum_polynomial}, choose a polynomial \(p\) such that
\[
  \forall z \in K,\qquad |p(z)-r(z)| < \eps/2 .
\]
Then for every \(z \in K\),
\[
  |p(z)-f(z)|
  \le |p(z)-r(z)| + |r(z)-f(z)|
  < \eps/2+\eps/2
  = \eps .
\]
Therefore
\[
  \exists p : \C[X],\quad
  \forall z \in K,\quad |p(z)-f(z)| < \eps .
\]
This proves the theorem.
\end{proof}

\section{Lean formalization notes}

\begin{itemize}
  \item The main theorem should use the existing Lean statement:
  \[
  \texttt{DifferentiableOn ℂ f U}
  \]
  as the holomorphicity assumption.

  \item The hardest formal ingredients are not the final \(\eps/2\) argument,
  but the analytic approximation lemmas:
  \[
  \texttt{cauchy\_integral\_approximated\_by\_pole\_sum}
  \]
  and
  \[
  \texttt{single\_pole\_polynomial\_approx}.
  \]

  \item If a shorter formalization route is available in mathlib, one can
  replace \cref{prop:rational_approximation} by an existing Runge-type rational
  approximation theorem, and then only formalize the connected-complement
  reduction from rational functions to polynomials.

  \item The final theorem depends conceptually on the following chain:
  \[
  \text{holomorphic on neighbourhood}
  \Rightarrow
  \text{finite pole-sum approximation}
  \Rightarrow
  \text{move poles to infinity}
  \Rightarrow
  \text{polynomial approximation}.
  \]
\end{itemize}

\end{document}